% GPSS complete-order boundary constructions and the necessary endpoint extension.

\subsection{Future limits and complete-order boundary constructions}
\label{sec:p09-gpss-complete-order-boundaries}

The extension from discrete complete orders needs formulas for nonattained
bounds. It also needs a change to the intermediate decomposition grammar:
the endpoint-free class in \cref{def:p09-gpss-decomposition} is not closed
under bounded relative complement on all complete orders. The counterexample
below has two actual endpoints. This section supplies future-only limit
formulas, an infimum-event compiler with a labelled actual lower endpoint,
and the bounded positive-chain result at its stated scope.
The modification announced after Proposition~2.4 concerns Kamp's
\(S,U\) result. It does not assert an arbitrary-complete-order
\(U\)-only future theorem. The local constructions below retain
their explicit language and endpoint restrictions
\cite[pp.~169--171]{GabbayPnueliShelahStavi1980}.

\begin{hblemma}{thm:p09-gpss-complete-order-bridge}{proved=annex:p9.gpss.complete.order.bridge}{The exact complete-order and temporal bridges}
The source condition that every nonempty bounded subset has an infimum and
supremum is equivalent to Kamp's \(\mathsf{Com}\) condition: every
nonempty lower-bounded subset has an infimum and every nonempty
upper-bounded subset has a supremum. The orders are nonempty and strict;
endpoints and adjacent points are permitted.

For formulas made from unary letters, Boolean operations and Kamp's
future constructor \(U(C,B)\), renaming its letters to the corresponding
GPSS letters and replacing that constructor by
\(B\mathbin{\mathsf U_s}C\) preserves truth on every linear order.
In particular the nested future word construction in
\cref{thm:p09-kamp-iii2-lemma1} supplies an \(L(U)\) formula when its
labels already have \(L(U)\) presentations. This bridge does not convert
an arbitrary elementary Kamp label, or a formula containing \(S\), into
a future-only label.
\end{hblemma}

\begin{hbdefinition}{def:p09-gpss-right-boundary-labels}{Nonempty right limits and right neighborhoods}
For a permitted future label \(A\), define the following actual \(L(U)\)
formulas:
\[
\begin{aligned}
 N^+A&=A\mathbin{\mathsf U_s}\top,\\
 \partial^+A&=\neg\bigl((\neg A)\mathbin{\mathsf U_s}\top\bigr),\\
 K^+A&=F\top\land\partial^+A.
\end{aligned}
\]
Here \(N^+A\) says that some strict right neighborhood has \(A\)
throughout its open interval. The symbol \(\partial^+\) agrees with
the unary future operation in \cref{def:p09-kamp-limit-operators};
\(K^+\) additionally requires a future point. These are abbreviations
for the displayed finite formulas, not new primitive connectives.
\end{hbdefinition}

\begin{hblemma}{thm:p09-gpss-right-limit}{proved=annex:p9.gpss.right.limit}{The precise right-limit semantics}
On every nonempty linear-order model,
\[
\begin{split}
 K^+A(x)\quad\Longleftrightarrow\quad&
 \exists b\,(x<b)\ \land\\
 &\forall b\bigl(x<b\longrightarrow
            \exists z\,(x<z<b\land A(z))\bigr).
\end{split}
\]
The predicate is tested at \(z\), with no condition on \(A(x)\).
Both \(\partial^+A\) and \(K^+A\) are false at a point with an
immediate successor. At a last point \(\partial^+A\) is true and
\(K^+A\) is false. For every actual \(x<y\), \(K^+A(x)\)
is equivalent to nonemptiness and coinitiality at \(x\) of the
\(A\)-points in \((x,y)\). This equivalence uses the whole-model
label \(A\), without restricting the quantifiers inside \(A\).
\end{hblemma}

\begin{hbdefinition}{def:p09-gpss-future-word-limits}{A future-only construction for finite-word accumulation}
A word is
\[
 W=(A_0,B_1,A_1,\ldots,B_n,A_n),\qquad n\geq0,
\]
with all point labels \(A_i\) and gap labels \(B_i\) in \(L(U)\).
Its realization is \(z_0<\cdots<z_n\), with \(A_i(z_i)\)
and \(B_i\) throughout \((z_{i-1},z_i)\). At \(n=0\)
there is just one labelled point. Define its starting formula by
\[
 E^+_{(A_0)}=A_0,\qquad
 E^+_W=A_0\land
             \bigl(B_1\mathbin{\mathsf U_s}E^+_{W^-}\bigr)
                    \quad(n\geq1),
\]
where \(W^-=(A_1,B_2,A_2,\ldots,B_n,A_n)\).
The formula \(E^+_W\) says that a realization starts at the current
point and places no upper bound on its end.

Define a second finite recursion:
\[
 \Gamma_{(A)}=K^+A,\qquad
 \Gamma_W=K^+(E^+_W)\land\Gamma_{W^-}\quad(n\geq1).
\]
Every recursive call removes the first point and first gap. Each resulting
formula uses only future until and Boolean operations on the supplied
labels; no past operator occurs.
\end{hbdefinition}

\begin{hblemma}{thm:p09-gpss-future-word-limit}{proved=annex:p9.gpss.future.word.limit}{Finite patterns arbitrarily close to the right}
For every word \(W\), \(\Gamma_W(x)\) is equivalent, over every
linear-order model, to
\[
 \exists d\,(x<d)\ \land\
 \forall b\left(x<b\to
   \exists z_0<\cdots<z_n
       \bigl(x<z_0\land z_n<b\land
                    (z_0,\ldots,z_n)\text{ realizes }W\bigr)\right).
\]
All realization conditions expand to the finite point and gap formula
in the definition.

Consequently Kamp's vacuous future limiting condition for the same
labelled word has the future-only presentation
\[
                       \neg F\top\ \lor\ \Gamma_W.
\]
For the original predicate-letter words this is exactly the semantic
condition of \(\operatorname{Lf}_n\) in
\cref{def:p09-kamp-limiting-classes}, including \(n=0\).
The proof supplies a different finite splice from the displayed
\(\Lambda_W^+\) recursion: that earlier recursion may contain \(S\)
inside a formula for a realization ending at the current point.
Here the formula is proved directly to remain in \(L(U)\).
\end{hblemma}

\begin{hblemma}{thm:p09-gpss-fixed-start-limit}{proved=annex:p9.gpss.fixed.start.limit}{Accumulation while retaining the initial point}
Let \(n\geq1\) and keep the word \(W\) above. The condition that
there is a later point and, for every \(b>x\), a realization of
\(W\) starts at \(x\) and ends strictly before \(b\), has the
future-only expression
\[
                      A_0\land N^+B_1\land\Gamma_{W^-}.
\]
The corresponding condition without the nonemptiness requirement has
expression
\[
        \neg F\top\ \lor\
                      (A_0\land N^+B_1\land\Gamma_{W^-}).
\]
Thus the fixed-start right-limit support used in Kamp's boundary
construction has an \(L(U)\) presentation whenever its labels do.
This statement retains the original initial point and the initial gap.
\end{hblemma}

\begin{hbdefinition}{def:p09-gpss-left-labelled-decomposition}{Decompositions with a labelled actual lower endpoint}
For actual \(x<y\), or actual \(x\) and formal \(+\infty\),
a \emph{left-labelled decomposition} is
\[
                           \alpha(x)\land D(x,y),
\]
where \(\alpha\) is a permitted future label and \(D\) is exactly a
decomposition of \cref{def:p09-gpss-decomposition}. The neutral
choice is \(\alpha=\top\). The same definition applies to a future
ray, with its permitted cofinal conditions. There is no label at a
formal point and no new individual parameter inside a label.

This is an explicitly enlarged intermediate grammar. The original
decomposition class has no point label at \(x\). All uses below
identify which of these two classes is the output class.
\end{hbdefinition}

\begin{hblemma}{thm:p09-gpss-left-labelled-closure}{proved=annex:p9.gpss.left.labelled.closure}{Positive assembly and future-ray complementation}
Finite disjunctions of left-labelled decompositions are closed under
finite conjunction and under the guarded existential concatenation of
\cref{thm:p09-gpss-positive-decomposition}, when the lower endpoint
is actual. At a concatenated midpoint \(z\), the right operand's
left label is conjoined with the required point label at \(z\).
All empty Boolean cases are retained.

On every linear order, every finite disjunction of left-labelled
future-ray decompositions has an \(L(U)\) expression, and its
complement has a left-labelled future-ray decomposition. Explicitly,
translate the disjunction to \(B(x)\) by
\cref{thm:p09-gpss-suffix-translation}, retaining every left label;
return
\[
              (\neg B)(x)\land D_{\top}(x,+\infty),
\]
where \(D_{\top}\) has no interior point, the single gap label
true, and no cofinal conjunct. Last points are included. The output
class here is the enlarged class just defined.
\end{hblemma}

\begin{hbdefinition}{def:p09-gpss-infimum-events}{Bounded infimum events and their exact statuses}
In a complete model let \(s<t\) be defined actual values, and put
\[
                     A_{s,t}=\{u:s<u<t,\ A(u)\}.
\]
The event \(\operatorname{Inf}_A(s,t)\) is undefined if a bound is
undefined, the bounds are not strictly ordered, or this set is empty.
Otherwise it returns \(c=\inf A_{s,t}\). It records one of three
statuses:
\[
\begin{array}{c|l}
 \mathrm{initial}&c=s,\\
 \mathrm{attained}&s<c<t,\quad A(c),\\
 \mathrm{unattained}&s<c<t,\quad\neg A(c).
\end{array}
\]
The initial status does not test \(A(s)\): \(s\) is excluded from
the occurrence set regardless of that truth value. Defined bounds and
values are domain points. Undefinedness is a separate symbol with no
place in their strict order.
\end{hbdefinition}

\begin{hblemma}{thm:p09-gpss-infimum-descriptors}{proved=annex:p9.gpss.infimum.descriptors}{All infimum statuses with future labels}
For actual \(s<t\), the following are exact descriptions of the
event and its status:
\[
\begin{array}{c|l}
 \mathrm{undefined}&\forall u\,(s<u<t\to\neg A(u)),\\
 \mathrm{initial}&c=s,\quad K^+A(s),\\
 \mathrm{attained}&s<c<t,\quad A(c),\
                    \forall u\,(s<u<c\to\neg A(u)),\\
 \mathrm{unattained}&s<c<t,\quad\neg A(c)\land K^+A(c),\
                    \forall u\,(s<u<c\to\neg A(u)).
\end{array}
\]
On a complete order exactly one row holds, with a unique value in
each defined row. Each row's equivalence to its specified infimum
property holds on any linear order on which that bound exists.
The point labels in these descriptions are actual \(L(U)\) formulas.
The initial row is precisely the new lower-endpoint label when \(s=x\).
An empty interval contributes only the undefined row.

For the unbounded future occurrence set \(\{u:s<u,\ A(u)\}\),
the same four statuses have the following formulas at \(s\),
in the same order:
\[
\begin{array}{c|l}
 \mathrm{undefined}&G\neg A,\\
 \mathrm{initial}&K^+A,\\
 \mathrm{attained}&(\neg A)\mathbin{\mathsf U_s}A,\\
 \mathrm{unattained}&(\neg A)\mathbin{\mathsf U_s}
                                      (\neg A\land K^+A).
\end{array}
\]
Over complete orders these are mutually exclusive and exhaustive,
including a last point. The interior boundary is the strict until
witness in each of the last two rows.
\end{hblemma}

\begin{hbtheorem}{thm:p09-gpss-infimum-case-compiler}{proved=annex:p9.gpss.infimum.case.compiler}{Finite infimum-event cases have explicit left-labelled decompositions}
Fix actual roots \(x<y\) and a finite acyclic list of events.
Each event is \(\operatorname{Inf}_A(s,t)\), where both arguments
are roots or earlier events, and \(A\) is a supplied future label.
One may also include tests \(\operatorname{Test}_A(s)\), returning
\(s\) if defined and satisfying \(A\), and undefined otherwise,
provided \(s\) is the root \(x\) or an earlier event, never \(y\).

Every Boolean condition on the definedness, the three infimum statuses,
and the equality/order of the defined values has an effective finite
disjunction of left-labelled decompositions equivalent to it on
complete models under \(x<y\).

The construction enumerates all definedness and status assignments and
all complete weak orders of the defined terms and roots. It retains
only the requested cases and compiles the displayed infimum descriptors
and test clauses. Equal values share one point. All labels at \(x\)
are conjoined into its left label. Every other nonroot value is
strictly between \(x,y\). Universal interval constraints are placed
on every merged gap and interior point they cover. No event symbol,
extra proposition letter, or endpoint parameter survives inside a label.
\end{hbtheorem}

\begin{hbdefinition}{def:p09-gpss-complete-positive-chain}{The complete-order positive-chain event recursion}
Fix a list of \(m\geq1\) required points
\(\Theta_1,\ldots,\Theta_m\) and preceding gaps
\(\Phi_0,\ldots,\Phi_{m-1}\), with no final-gap condition.
Let \(V_{\Theta;\Phi}(a,b)\) be the following Boolean event
expression. For the empty point list it is true. For a nonempty
list and undefined or non-strictly ordered bounds it is false.
Otherwise form the point word
\[
                     W=(\Theta_1,\Phi_1,\Theta_2,\ldots,
                                      \Phi_{m-1},\Theta_m)
\]
and the events
\[
\begin{aligned}
 c&=\operatorname{Inf}_{E^+_W}(a,b),&
 h&=\operatorname{Inf}_{\neg\Phi_0}(a,c),\\
 j&=\operatorname{Test}_{N^+\Phi_0}(c),&
 k&=\operatorname{Test}_{\Phi_0\land N^+\Phi_0}(c).
\end{aligned}
\]
For \(m=1\), \(W=(\Theta_1)\). For \(m>1\) let
\[
 R=V_{(\Theta_2,\ldots,\Theta_m);
                  (\top,\Phi_2,\ldots,\Phi_{m-1})}(c,b);
\]
for \(m=1\) let \(R=\mathsf{true}\). Thus the shortened
list has a true initial gap and retains every gap between
its remaining points. Put
\[
\begin{split}
 V_{\Theta;\Phi}(a,b)={}&
 \operatorname{Def}(c)\land\neg\operatorname{Def}(h)\land R\\
 &{}\land\bigl[
       \operatorname{Att}(c)
       \lor(\operatorname{Init}(c)\land\operatorname{Def}(j))
       \lor(\operatorname{Unatt}(c)\land\operatorname{Def}(k))
                    \bigr].
\end{split}
\]
Here \(\operatorname{Att},\operatorname{Init},\operatorname{Unatt}\)
are the recorded statuses of this particular infimum event, false
when it is undefined. The recurrence removes one point per call,
even if an initial event returns the same lower endpoint.
All symbols in this display are metalinguistic event expressions
to be compiled by \cref{thm:p09-gpss-infimum-case-compiler}.
\end{hbdefinition}

\begin{hbtheorem}{thm:p09-gpss-complete-positive-chain}{proved=annex:p9.gpss.complete.positive.chain}{A determining construction for bounded positive chains on complete orders}
On actual \(a<b\) in a complete model, the event expression just
defined holds exactly when there are points
\[
                           a<z_1<\cdots<z_m<b
\]
with \(\Theta_i(z_i)\), with \(\Phi_0\) throughout
\((a,z_1)\), and with \(\Phi_i\) throughout
\((z_i,z_{i+1})\) for \(1\leq i<m\).
The empty point list imposes no condition.

For \(m\) points the event list contains at most \(4m\)
nonroot nodes in addition to its two roots. It uses only
future labels, infima and tests. Its Boolean decision
therefore has an explicitly constructible finite
disjunction of left-labelled decompositions, as does
its relative complement under \(a<b\).

The attained infimum supplies a first point. An initial
or interior unattained infimum instead supplies points
arbitrarily close on its right; the neighborhood and
interior point tests in the displayed recurrence preserve
the whole initial gap. The suffix can be realized before
the upper bound independently of the provisional global
future realization, and the proof supplies the exact
splice joining them. This replaces the finite-interval
first-occurrence assumption of
\cref{thm:p09-gpss-canonical-prefix} for this positive-chain
component. A final-gap requirement remains separate.
\end{hbtheorem}

\begin{hbtheorem}{thm:p09-gpss-complete-bounded-obstruction}{proved=annex:p9.gpss.complete.bounded.obstruction}{Why the original bounded output class is insufficient}
The original endpoint-free decomposition class is not closed under
relative complement on arbitrary complete orders, even with actual
bounded endpoints. It already fails for
\[
 D_A(x,y)=x<y\land
       \exists z\bigl(x<z<y\land
                \forall u(x<u<z\to A(u))\bigr).
\]
This is an original decomposition with one true-labelled interior
point, initial gap \(A\), and final gap true.

Take the complete order
\[
                  T=\{0,2\}\cup\{1/n:n\geq1\},
                  \qquad x=0,\quad y=2.
\]
Vary \(q_0(1/n)\) independently, take \(A=q_0\), fix \(q_0\)
at the endpoints, and fix all other predicates throughout \(T\).
No finite disjunction of original decompositions agrees with
\(x<y\land\neg D_A(x,y)\) for every such interpretation.
This is a failure at an actual lower endpoint with no immediate
successor, distinct from the formal-left-endpoint obstruction in
\cref{thm:p09-gpss-prefix-complement-obstruction}.
\end{hbtheorem}

\begin{hblemma}{thm:p09-gpss-initial-gap-complement}{proved=annex:p9.gpss.initial.gap.complement}{The complete initial-neighborhood complement with its endpoint label}
On every linear-order model, for actual \(x<y\),
\[
 \neg D_A(x,y)\quad\Longleftrightarrow\quad
       \operatorname{Adj}(x,y)\ \lor\ K^+(\neg A)(x),
\]
where \(\operatorname{Adj}(x,y)\) means
\(x<y\land\neg\exists z(x<z<y)\).
Thus its guarded complement is the disjunction of two explicitly
left-labelled decompositions: the zero-point word with false gap and
true left label; and the zero-point word with true gap and left label
\(K^+(\neg A)\). These cases include adjacency, a nonattained
initial limit, and every truth value of \(A(x)\) and \(A(y)\).
\end{hblemma}

\begin{hbexample}{def:p09-gpss-boundary-examples}{Attained and unattained beginnings}
On \(\mathbb R\), with bounds \(0<2\), take successively
\(A=\varnothing\), \(A=(0,1)\), \(A=\{1\}\), and
\(A=(1,3/2)\). The infimum-event statuses are respectively
undefined, initial, attained at 1, and unattained at 1. Replacing
\(A(0)\) alone changes none of these statuses. On \(0<1\) with
no intermediate point the event is undefined for every interpretation.

For \(W=(A,B,C)\), the word-limit formula is
\[
       K^+\bigl(A\land(B\mathbin{\mathsf U_s}C)\bigr)
                                      \land K^+C.
\]
On the real line with \(A=(0,1)\), \(B=\mathbb R\), and
\(C=\{2\}\), its first conjunct holds at zero but its second
does not; the required two-point patterns cannot occur arbitrarily
close to zero. With \(A=C=(0,1)\) both conjuncts hold and every
right neighborhood contains a realizing pair.

For an event-list example take
\(i=\operatorname{Inf}_P(x,y)\) and
\(j=\operatorname{Inf}_Q(i,y)\).
The case that \(i\) is initial and \(j\) is attained compiles to
\[
 K^+P(x)\land x<y\land
 \exists c\bigl(x<c<y\land Q(c)\land
                   \forall u(x<u<c\to\neg Q(u))\bigr).
\]
Its left label is \(K^+P\), its one point label is \(Q\),
and its two gaps are \(\neg Q,\top\). The infimum descriptors
identify \(i=x\) and \(j=c\) in both directions.

In the positive-chain recurrence with one point, bounds \(0<2\),
and \(\Theta_1=(1,3/2)\) on \(\mathbb R\), the candidate
infimum is the unattained interior point \(c=1\). With
\(\Phi_0=\mathbb R\) the chain exists. With
\(\Phi_0=\mathbb R\setminus\{1\}\) it does not: every
possible first point lies above 1, so its initial gap contains
the failed point 1. The interior test in \(k\) distinguishes
these interpretations even though both open side intervals
can satisfy their gap conditions.
\end{hbexample}

\begin{proofboundary}{def:p09-gpss-complete-boundary-continuation}{What the boundary constructions establish}
The infimum compiler translates a supplied finite list; it does not
by itself construct a list deciding every decomposition.
\Cref{thm:p09-gpss-complete-positive-chain} supplies such a list
for the whole positive chain without its final gap.
A general final gap involves upper-boundary information that is
not covered by those finite-list premises. Consequently the proved
local statements do not license a full regional negation induction
in this future-only grammar.

For a nonempty occurrence set in \((a,b)\), an unattained supremum
\(c\) is either the outer endpoint \(b\) or an interior point.
Its left approach requires
\[
 \forall r<c\ \exists u\,
          (a<u<c\land r<u\land A(u)).
\]
This includes a real witness by taking \(r=a\). An interior
supremum also requires \(\neg A(c)\) and absence of occurrences
in \((c,b)\); an attained interior supremum instead requires
\(A(c)\) and that same absence. These are first-order boundary
conditions, not admitted future labels at \(c\).
In particular left accumulation at \(c\) cannot itself be a
future-only label there: on \(\mathbb R\), interpretations
\(A=(-1,0)\) and \(A=\varnothing\) agree at and after zero
and disagree on left accumulation at zero.

Kamp's complete-order construction permits past-dependent elementary
labels at such boundaries, and proves completeness for the
\(S,U\) basis. That valid result is not a promise that the past
operators can be removed from every complete-order future presentation.
The countertheorem above also prevents an inference that the original
endpoint-free GPSS Main Lemma extends unchanged. The source's announced
complete-order modification and any separately stated future-language
comparison must retain their actual operator bases
\cite[p.~171]{GabbayPnueliShelahStavi1980}.
\end{proofboundary}
